\section{Simulator}

In order to be able to develop an LLM scheduler, we need access to many high-performance GPUs to test our schedule configurations.
Unfortunately, GPU-hours are prohibitively expensive and using such hours for ``dummy'' workloads, to test a scheduler, would be considered a waste of resources. %CITE
Instead, we utilise the latest simulator technology to estimate the performance of LLM inference given a set of input parameters.
Several simulators already exist that attempt to model the performance of an LLM request by using estimation techniques.

% Previous work has also used schedulers too (reference the work here).

\textbf{Vidur} \autocite{agrawalVidurLargeScaleSimulation2024} is an LLM simulator developed by Microsoft, designed to address the high cost and complexity of experimentally finding the best deployment configuration of LLMs.
This is achieved by profiling a model-gpu combination for expected throughput on given operations (different prefill and decode requests).
Then, using a small machine learning model, the system predicts the likely performance with high degrees of accuracy for different parrallelism configurations.

\textbf{Shallowsim} \autocite{icezack12Icezack12Shallowsim2025} is a specialist simulator designed to estimate the compute times for different LLM request types.
Specifically built for DeepSeek-V3/R1 models, it has native support for simulating MoE architectures.
Unlike Vidur, Shallowsim requires no profiling phase, rather a list of device performance metrics is provided e.g. \texttt{fp8} performance.
These metrics are then used to compute the expected performance for the prefill and decode phase separately

We select Shallowsim for its transparency of calculations (simplicity), as it uses no machine learning to predict performance.
Instead, simple mathematical operations are used, which allows us to inspect performance and modify the calculations as necessary.
The Shallowsim simulator also already has native support for Mixture of Expert models, which is our key novel design goal for the scheduler.

%And requires complex and expensive profiling

\subsection{Shallowsim Implementation}

%\todo{
%\textit{This is just describing how shallowsim does its maths, will finish very last as it is not strictly necessary for the report - Might even include in the appendix if too long.}
%}

%\textbf{Prefill}
%
%Dense-MLA (dense Multi-head/Multi-query Attention compute),
%TP-MLA (Dense-MLA plus the Tensor-Parallel all-reduce),
%Dense-MLP (feed-forward network compute),
%Shared-Expert (MoE expert run that every token visits),
%Routed-Expert (MoE expert run that only selected tokens visit),
%Every token first passes through the Shared Expert, but a subset of tokens are also routed to extra specialists. This is modelled with the Routed Expert function.
%%Dispatch (All-to-All that ships tokens out to remote experts),
%%Combine (All-to-All that brings tokens back),
%%Sending activations to the respective routed expert nodes is modelled using the Dispatch step, the Combine step then models how processed activations are collated back together. These two functions model the network overhead of LLM queries.
%
%\textit{Compute Time
%-	Dense MLA, Dense MLP (Dense Layers)
%-	TP-MLA, Shared Expert, Routed Expert (Sparse Layers)
%Communication Timemake 
%-	Dispatch, Combine (Sparse Layers)
%Overlap Penalties
%-	If Combine is longer than TP-MLA + Shared Expert
%-	If Dispatch is longer than Routed Expert
%Sum
%-	Compute + Communication + Any Overlap = Prefill Time
%}
%
%\subsubsection{Decode}
%
%Load-KV,
%Dense-MLA (as prefill),
%Sparse-MLA,
%Dense-MLP (as prefill),
%Shared-Expert (as prefill),
%Routed-Expert (as prefill),
%Dispatch (as prefill),
%Combine (as prefill),
%
%\textit{Compute Time
%-	Dense MLA, Dense MLP (Dense Layers)
%-	Sparse MLA, Shared Expert, Routed Expert (Sparse Layers)
%Communication Time
%-	Dispatch, Combine (Sparse Layers)
%Overlap Penalties
%-	TODO
%Sum
%-	TODO = Per token decode time 
%}

\subsection{Modifactions}

We make several small modifications to the Shallowsim simulator to support our use-case of efficiently profiling system configurations.
As-built Shallowsim is designed to profile several different parallelism configurations and batch sizes, then report back the overall fastest configuration.
We need support for profiling a specific parallelism configuration and batch size, as we may not always be able to use the most optimal for a given exact sequence length, but rather a range of sequence lengths.

To do this, we extract all the mathematical operations present in the existing code and form procedures with specific input parameters, e.g., $\texttt{prefill\_a2a}(\texttt{gpu}, \texttt{seq\_len}, \texttt{tp\_num})$.
Here, we precisely specify the gpu model, input sequence length, and tensor parallelism degree.

Given the modifications, shallowsim now supports the following functions:
\begin{itemize}
    \item \textbf{Prefill Performance}\ \texttt{-> Milliseconds}\\
    \texttt{prefill\_time(gpu, seq\_len, kv\_cache\_rate, tp\_num, dp\_num)}
    
    \item \textbf{Decode Performance}\ \texttt{-> Milliseconds}\\
    \texttt{decode\_time(gpu, tp\_num, bs\_num, seq\_len, decode\_len, moe\_num, device\_num)}
    
    \item \textbf{Maximum Batch Size}\  \texttt{-> Int}\\
    \texttt{check\_max\_bs(gpu, tp\_num, bs\_num, seq\_len, decode\_len, moe\_num, device\_num)}
\end{itemize}

\subsection{Usage}
\label{chap:simulator-usage}

\subsubsection{Prefill Phase}
Using the Shallowsim simulator we can measure the time $t$ (in milliseconds) to compute a single prefill request for a given set of parameters, and convert this to throughput R (in requests per second) via
$R_{\text{prefill}} = \frac{1000}{t}$.

% include sample parameters for prefill function

\subsubsection{Decode Phase}
Shallowsim has support for calculating the time in milliseconds for a single token generation step.
As we would typically expect several tokens to be generated, we cannot directly compute requests per second using Shallowsim, instead we must perform additional computations.
The decode time of a single token is affected by the input/current sequence length, therefore we must measure the single token compute time for every token to be generated in a request.

% graph of sequence length affecting decode time

By measuring $t_i$ the compute time (in milliseconds) for the $ith$ token, the total decode time for a request generating N tokens is $t_{\text{decode}} \;=\;\sum_{i=1}^{N} t_i\,$.
We can then convert this total time into throughput $R_{\text{decode}}$ (in requests per second), using the same method as prefill.
$R_{\text{decode}}
\;=\;\frac{1000}{t_{\text{decode}}}
\;=\;\frac{1000}{\sum_{i=1}^{N} t_i}\,$,
where $R_{\text{decode}}$ represents how many full decode requests per second the configuration can sustain.

